\documentclass[11pt]{article}
\usepackage{amsmath,amssymb,amsthm,geometry}
\geometry{margin=1in}\newtheorem{proposition}{Proposition}
\title{Weak passage with an added inactive strong cohort}
\author{AN04 closure campaign}\date{Version 1}
\begin{document}\maketitle
\paragraph{Leading system, Route A.}
Fix a point of WIT's witness rectangle. Write $(X,Y,U,V)$ for
the coherent atom coordinates on its logistic clock, $n$ for weak mass,
$P=\Phi(\rho X)$, $Q=\Phi(U)$, $K_x=K(\rho X)$, $K_u=K(U)$,
$A=\zeta-X$, $Z=X-Y$, $L=Y+nK_u$, $o=\rho(1-n)K_x/2$.
Take $s_0=\inf J-h$ with $0<h\le1$ small enough that $X<\zeta$
on $I=[s_0,\sup J]$; such $h$ exist by WIT's strict compact margin.
Add mass $\varepsilon$ at $(x_b,y_b)$ at $s_0$ to the unchanged
unit core, retaining the same weak atom and seed $n(s_0)$.
This literal addition makes $M(s_0)=1+\varepsilon$.
A replacement of core mass would give a different derivative.
All perturbations below are right derivatives at zero mass.

\paragraph{Passive trajectory and an exterior range.}
At zero added mass its label $(b_x,b_y)$ solves
\begin{align*}
 2b_x'&=\rho\phi(\rho b_x)-\rho F(\rho b_x,\rho X)
      +\lambda[nV+(1-n)b_y]\Phi(\rho b_x),\\
 2b_y'&=-L+\rho(1-n)K(\rho b_x).
\end{align*}
Set $D=|I|$, $B=1+\sup_I (|L|+\rho K_x+\lambda |V|)$.
For $y_b\ge0$ one has $b_y\ge-BD/2$ and
$b_x'\ge-B/2-\lambda BD/4$, by $0\le F\le K_x$.
In particular the explicit range
\begin{equation}\label{an04swpv1:outside}
 y_b\ge0,\qquad x_b>\zeta+1+BD/2+\lambda BD^2/4
\end{equation}
stays strictly inactive throughout $I$. These are integrated scalar
bounds in the reached coherent field. All coordinates are finite on
$I$ by linear growth. The finite-atom WIT field is continuously
differentiable: $F_t(t,r)=(r-t)_+\phi(t)$ and
$F_r(t,r)=\Phi(\min(t,r))$, also at equality. Standard finite-time
continuous dependence keeps the core active and bridge inactive for
all sufficiently small $\varepsilon$ for each fixed pair in this range.

\begin{proposition}[Leading orbit response]\label{an04swpv1:response}
Let $(j,k,l,m)$ be the variations of the core and weak coordinates,
respectively, in their displayed order. Their initial data vanish.
Put $d=e^{-(s-s_0)}$, $c=d-1$. Then
\[
 Y_1=k+cY+b_y,\quad K_1=\rho Pj+cK_x+K(\rho b_x),\quad
 S_1=\rho Pj+cK_x+F(\rho X,\rho b_x).
\]
With $p_x=\phi(\rho X)$ and $p_u=\phi(U)$ the exact linear response is
\begin{align}
 2j'={}&dX-\rho^3 X p_x j-\rho S_1
  +\lambda P[nm+(1-n)k]+\lambda[nV+(1-n)Y]\rho p_xj,\nonumber\\
 2k'={}&dY-Y_1-nQl+\lambda(1-n)Pj,\nonumber\\
 2l'={}&-Up_u l-nQl-(Y_1-dV)Q-Yp_u l-\lambda(1-n)l,\nonumber\\
 2m'={}&\rho K_1-\lambda m+dK_u.\label{an04swpv1:linear}
\end{align}
\end{proposition}
\begin{proof}
$M'=M(1-M)$ and the normalized strong weights are fixed, so the
derivatives of total, core and bridge mass are $d,c,1$.
Weak mass has zero derivative. Differentiate WIT popsystem at the
core and weak receivers. The core self overlap is $K(\rho X)$,
including receiver and donor motion; the other strong donor enters
with coefficient one. Hence $Y_1,K_1,S_1$ are the full moment and
receiver-overlap variations. The weak self overlap is $nK(U)$.
In particular at injection the bridge supplies $-y_b/2$ through
the output moment, $-\rho[\phi(\rho X)+\rho x_bP]/2$ through input
overlap, and $+y_b$ through full-output lag regression. The two
strong velocity equations also gain $X/2,Y/2$ from the mass defect.
More precisely the direct increment per unit added mass in the
finite-noise static regression $r_{21}/q$ is, by P stein and Lexcharts,
\[
 \frac{\sin(qy_b)}{1+q^2}
 \left[K(\cos(qx_b)/q)+q\cos(qx_b)\Phi(\cos(qx_b)/q)\right].
\]
Here the receiver is specified in the strong chart at fixed finite
$(x_b,y_b)$. Gaussian integration by parts on the full regression,
or the two cited identities, gives this expression exactly.
It tends to $y_b$ as $q\downarrow0$. The finite-$q$ literal
coefficient need not equal $y_b$. This static identity does not transfer
the perturbed trajectory to training.
\end{proof}

\paragraph{Leading rates.}
Write $a_1=cA-j$, $L_1=Y_1+nQl$,
$o_1=\lambda(1-n)Pj/2$ and
$C_0=(A-\zeta)A+YA/2-ZX/2$. On $J$, the exact variations are
\begin{align}
 D_1&=-(a_1 L+A L_1)/2-d C_0,\nonumber\\
 D_2&=c(ZX'+Ao)+(-a_1-k)X'
       +Z(j'-dX/2)-jo+Ao_1,\nonumber\\
 D&=D_1+D_2=\partial_\varepsilon\mathcal E'_\zeta|_{0+}.
 \label{an04swpv1:rates}
\end{align}
Indeed WIT generalrates gives $d_1=-A_\zeta L/2+(1-M)\mathcal C$.
For source two, only the active core is integrated, with weight
$M_c$, input coefficient $\zeta-A_\zeta-y_c$, and output coefficient
$\zeta-x_c$. Its source-two input variation is $j'-dX/2$.
This derivation differentiates the rate integrals with their actual
weights and current gates. Outside the separated range, at any time with $b_x\ne\zeta$ and
strict core activity, add $A_b=(\zeta-b_x)_+$ to $a_1$ everywhere.
Also add $\mathbf 1_{b_x<\zeta}(\zeta-A-b_y)b_x'+A_b
\rho(1-n)K(\rho b_x)/2$ to $D_2$.
At zero mass $M=1$, so $b_x'$ is entirely source two. Strict
separation in the stated range avoids all gate-contact derivatives.

\paragraph{A signed subset and the window test.}
At injection, extending the active formulas to $s_0$,
\begin{equation}\label{an04swpv1:initial}
 D(s_0)=-Ay_b/2-C_0-Z\rho[\phi(\rho X)+\rho x_bP]/2 .
\end{equation}
Here $Z>0$ by the preceding AN01 result. To give a signed result
\emph{inside} the window put $s_L=\inf J$. Compute $B,D$ in the
exterior bound on the larger interval $[s_L-1,\sup J]$ (the bound only
uses the donor coordinates). Choose $y_b\ge0$ and $x_b$ satisfying
that enlarged bound and
$x_b>2(|C_0(s_L)|+1)/(\lambda Z(s_L)P(s_L))$.
At $h=0$ the right side of \eqref{an04swpv1:initial} is below $-1$.
Continuous dependence of the passive and inhomogeneous linear initial
value problems on their initial time gives, for each fixed pair,
$h_0>0$ and $\eta>0$ with $D<0$ on
$[s_L,s_L+\eta]\subset J_-$ for $0<h<h_0$.
The negative change strengthens descent on this subinterval.
For the entire $J_-$ or $J_+$ the unpaid respective tests are
\begin{equation}\label{an04swpv1:missing}
 \sup_{s\in J_-}D(s)<0,\qquad \inf_{s\in J_+}D(s)>0 .
\end{equation}
Here $D$ is fully specified by \eqref{an04swpv1:linear}--\eqref{an04swpv1:rates}.
For example the $-A nQl/2$ contribution competes with the direct
$-A b_y/2$, while $l'$ contains $-QY_1/2$ and $m'$ contains
$+\rho K_1/2$. No same-population order theorem signs this forced
system. For any fixed separated pair, compactness and WIT's strict
margins do preserve $d_1<0<d_2$ on $J$ and the opposite total-rate
signs on $J_-,J_+$ for small $\varepsilon$. These are zeroth-order
signs of the perturbed rates, distinct from \eqref{an04swpv1:missing}.
For the rebound between any two specified times $t_-,t_+$ the exact
first-order increase test is $\int_{t_-}^{t_+}D(s)\,ds>0$.
An optimized minimum or maximum additionally needs its optimizing
times treated; the supplied window margins do not identify them.

\paragraph{Leading system, Route B.}\label{an04swpv1:full}
Prescribe the entire RMS join, $M=1$, and weak mass $0<N_0<1/2$
at $(u_0,v_0)=(0,0)$ as a concrete leading seed.
The logistic half delay is exactly
$\lambda^{-1}\log[(1-N_0)/N_0]$. At entry $K_w=N_0\phi(0)$ and $V_w=0$;
on the aligned strong gate, with $t=\rho\zeta$, the doubled velocity is
$\rho[K(t)-\langle F(t,\rho x)\rangle_b]-\lambda N_0\zeta\Phi(t)$.
This follows directly from WIT, including the weak increment.
On a characteristic interval the weak atom remains an atom, with
$2u'=\phi(u)-NK(u)-Y_s\Phi(u)-\lambda(1-N)u$,
$2v'=\rho K_s-\lambda v$ and
$2Y_s'=-Y_s-NK(u)+\rho(1-N)K_s$.
A finite-amplitude increase of a specified rebound needs
$[\mathcal E_{\rm split}-\mathcal E_{\rm coh}]_{t_-}^{t_+}>0$
with matched clocks, where
$\mathcal E=A_\zeta^2/2-\zeta A_\zeta+\int y(\zeta-x)_+d\mu_s$.
The first three moment equations leave both exposure integrals unpaid:
their derivatives involve the active-rank input velocities.
The gate sign at injection, and rank return when $N$ is identically zero,
give no sign for this endpoint difference. No weak-seed scaling
or trained-flow depth conclusion follows.
\end{document}
